% TEMPLATE-MILC-v3.tex - plantilla maestra UNICA de las guias MILC v3.
%
% La usan las cuatro familias de guias, cada una con su driver:
%   · Grados 6-11          -> scripts/build-guias-g11.py
%   · Bebras               -> scripts/build-guias-bebras.py
%   · Semillero            -> scripts/build-guias-semillero.py
%   · Territorio interior  -> scripts/build-guias-territorio-interior.py
%
% Antes habia tres copias casi identicas (~400 lineas iguales, 6 distintas):
% cada mejora habia que aplicarla tres veces, y en la practica no se hacia
% (el motor de recursos vivia solo en dos de ellas). Lo que varia entre
% familias esta parametrizado en 6 placeholders que cada driver llena
% (se nombran aqui SIN su sintaxis real: el builder reemplaza por texto plano,
% tambien dentro de los comentarios, y un valor multilinea rompe el preambulo):
%
%   HEADER_IZQ        encabezado izquierdo de pagina
%   HEADER_DER        encabezado derecho de pagina
%   PORTADA_ETIQUETA  rotulo sobre el numero grande (GRADO/MOMENTO/MODULO)
%   PORTADA_SERIE     linea de serie de la portada
%   PORTADA_ID        identificador de la guia en la portada
%   PORTADA_CIRCULO   el circulito de grado (°), o vacio si no aplica
%   PDF_TITULO        titulo en texto plano (sin \textbf ni comillas TeX) para
%                     los metadatos del PDF (/Title); lo produce lib_xelatex.texto_pdf
%
% Composicion (auditoria 2026-09-03): las bandas de estacion piden espacio
% (\Needspace*) y prohiben el salto justo despues (\nopagebreak), asi no quedan
% huerfanas al pie; las cajas largas (softbox, drawbox, infoband, puentebox) son
% `breakable`, asi una caja de media pagina ya no arrastra todo a la siguiente
% dejando la anterior al 40 %. Todo texto sobre mostaza o verde va en milcNegro
% (blanco sobre #E5B400 da 1,93:1 y sobre #58A923 2,95:1; ninguno pasa WCAG AA).
% El resto de claves las documenta content/guias/_SCHEMA.md. El script valida
% que no queden placeholders sin reemplazar antes de escribir el archivo final.

% Etiquetado PDF (latex-lab): /Lang, estructura y texto alternativo de las
% figuras, para lectores de pantalla. Debe ir ANTES de \documentclass.
\DocumentMetadata{lang=es-CO,tagging=on}
\documentclass[11pt,letterpaper]{article}
% headheight=24pt: la cabecera derecha lleva dos lineas (identidad de la guia
% y estacion actual). headsep se acorta en la misma medida para que el bloque
% de texto quede exactamente donde estaba con headheight=12pt.
\usepackage[margin=2.0cm,headheight=24pt,headsep=13pt]{geometry}
\usepackage[table]{xcolor}
\usepackage{fontspec}
\usepackage[spanish]{babel}
\usepackage{tikz}
% Librerías TikZ para los diagramas de las guías (bloque `recursos:`):
%  · babel  → babel en español vuelve activos caracteres como «>», y TikZ los usa
%             en sus opciones (>=stealth, ->). Sin esta librería, cualquier
%             diagrama con flechas revienta con «\language@active@arg> has an
%             extra }». Debe ir ANTES de que se dibuje nada.
%  · positioning        → sintaxis «below left=2pt and 3pt».
%  · decorations.pathreplacing → llaves ({brace}) para señalar rangos.
\usetikzlibrary{babel,positioning,decorations.pathreplacing}
\usepackage{graphicx}
% Circuitos: el trabajo de electrónica se hace hoy con micro:bit y sus sensores
% integrados (MakeCode, sin cableado), así que no se necesitan esquemáticos.
% Si algún día entran montajes con protoboard, basta descomentar esta línea:
% los diagramas .tex/.tikz declarados en `recursos:` ya se incrustan solos.
% \usepackage{circuitikz}
\usepackage[most]{tcolorbox}
\usepackage{tabularx}
\usepackage{array}
\usepackage{fancyhdr}
\usepackage{titlesec}
\usepackage[document]{ragged2e}  % bandera a la derecha en todo el cuerpo
\usepackage{enumitem}
\usepackage{microtype}
\usepackage{etoolbox}
\usepackage{needspace}
% hyperref al final del preambulo: metadatos del PDF (titulo, autor, idioma,
% familia), marcadores por estacion (\pdfbookmark) y enlaces sin recuadro.
\usepackage[hidelinks,
  pdftitle={Gráficos — la imagen de los datos},
  pdfauthor={PhD. Álvaro Cárdenas Orozco},
  pdfsubject={Tecnología e Informática · Grado 9 · Guía 26 · MILC v3},
  pdflang={es-CO},
  bookmarksopen=true,bookmarksnumbered=false]{hyperref}

% Helvetica Neue no trae la flecha U+2192, asi que cada «→» escrito literalmente
% en el YAML se compilaba a NADA, en silencio (462 flechas en 61 guias: rutas del
% tipo «paleta → tipografia → lockup» salian rotas). Mapeamos el caracter a su
% version matematica, que si tiene glifo. Asi el autor puede seguir escribiendo
% «→» comodamente en el YAML.
\usepackage{newunicodechar}
\newunicodechar{→}{\ensuremath{\rightarrow}}
% Mismo problema con los simbolos de marca y vinetas: el «✓» de «una palomita ✓»
% salia como cuadrito vacio en el PDF de todas las guias que lo usaban.
\usepackage{pifont}
\newunicodechar{✓}{\ding{51}}
\newunicodechar{✗}{\ding{55}}
\newunicodechar{✔}{\ding{52}}
\newunicodechar{•}{\textbullet}
\newunicodechar{≥}{\ensuremath{\geq}}
\newunicodechar{≤}{\ensuremath{\leq}}

\setmainfont{Helvetica Neue}
\setsansfont{Helvetica Neue}
\setlength{\parindent}{0pt}
\setlength{\parskip}{6pt}
% Sin particion de palabras: en 6.º-7.º una palabra cortada por guion al final
% de linea («Comuni-cacion») frena la lectura mucho mas que una linea corta.
\hyphenpenalty=10000
\exhyphenpenalty=10000
\pretolerance=2000
\tolerance=3000
\setlength{\emergencystretch}{3em}
\linespread{1.10}
\renewcommand{\arraystretch}{1.25}
\setlist{leftmargin=5mm,itemsep=1mm,topsep=1mm}
\newcolumntype{Y}{>{\RaggedRight\arraybackslash}X}

\definecolor{milcMagenta}{HTML}{E6007E}
\definecolor{milcVino}{HTML}{5A0038}
\definecolor{milcTurquesa}{HTML}{0093A5}
\definecolor{milcVerde}{HTML}{58A923}
\definecolor{milcMostaza}{HTML}{E5B400}
\definecolor{milcOcre}{HTML}{B86600}
\definecolor{milcNegro}{HTML}{241816}
\definecolor{milcGris}{HTML}{F7F7F4}
\definecolor{milcLinea}{HTML}{E8E2DD}
\definecolor{milcGradePrimary}{HTML}{7C3AED}
\definecolor{milcGradeDark}{HTML}{4C1D95}
\definecolor{milcGradeSoft}{HTML}{EDE9FE}
\definecolor{milcGradeAccent}{HTML}{CFE7FF}
\definecolor{milcGradeText}{HTML}{FFFFFF}
\color{milcNegro}

\pagestyle{fancy}
\fancyhf{}
% Cabecera de dos lineas: identidad de la guia arriba y, a la derecha y en
% gris, la estacion en curso (\markright desde \milcestacion). Antes de la
% primera estacion la segunda linea va vacia; el \strut mantiene la altura
% para que la linea de arriba no baile entre paginas.
\lhead{\small Tecnología e Informática · Grado 9\\[-1pt]{\footnotesize\strut}}
\rhead{\small Guía 26 · MILC v3\\[-1pt]{\footnotesize\strut\color{milcNegro!65}\rightmark}}
\cfoot{\small\thepage}
\renewcommand{\headrulewidth}{0pt}

% Sobre mostaza (#E5B400) y verde (#58A923) el texto va en milcNegro; sobre el
% resto de la paleta, en blanco. Se usa dentro de fonttitle/fontupper, donde un
% \color posterior manda sobre coltitle.
\newcommand{\milctintasobre}[1]{%
  \ifboolexpr{test{\ifstrequal{#1}{milcMostaza}} or test{\ifstrequal{#1}{milcVerde}}}%
    {\color{milcNegro}}{\color{white}}}

\newtcolorbox{titlebox}[1]{
  enhanced,colback=#1,colframe=#1,arc=7mm,boxrule=0pt,
  left=7mm,right=7mm,top=3mm,bottom=3mm,
  before skip=8pt,after skip=8pt,
  fontupper=\sffamily\bfseries\Large\color{white}
}

\newtcolorbox{softbox}[3]{
  enhanced,breakable,pad at break=3mm,
  colback=#2,colframe=#1,borderline west={4pt}{0pt}{#1},
  arc=2mm,boxrule=.4pt,left=4mm,right=4mm,top=3mm,bottom=3mm,
  before skip=6pt,after skip=8pt,title={#3},coltitle=white,
  colbacktitle=#1,fonttitle=\sffamily\bfseries\milctintasobre{#1},fontupper=\RaggedRight,
  attach boxed title to top left={xshift=3mm,yshift=-2mm},
  boxed title style={arc=2mm, boxrule=0pt}
}

\newtcolorbox{drawbox}[2]{
  enhanced,breakable,pad at break=4mm,
  colback=white,colframe=#1,arc=4mm,boxrule=1pt,
  left=5mm,right=5mm,top=4mm,bottom=4mm,before skip=8pt,after skip=8pt,
  title={#2},coltitle=white,colbacktitle=#1,fonttitle=\sffamily\bfseries\milctintasobre{#1},
  fontupper=\RaggedRight,
  attach boxed title to top left={xshift=4mm,yshift=-2mm},
  boxed title style={arc=3mm, boxrule=0pt}
}

\newtcolorbox{infoband}[1]{
  enhanced,breakable,pad at break=2mm,colback=#1!8,colframe=#1!50,arc=2mm,boxrule=.5pt,
  left=4mm,right=4mm,top=2mm,bottom=2mm,before skip=4pt,after skip=4pt,
  fontupper=\small\RaggedRight
}

% Puente narrativo entre secciones (contrato editorial Conéctate).
% Da continuidad: "Acabas de X. Ahora vas a Y porque Z."
% Visualmente: banda izquierda en color del grado, texto en itálica pequeña.
\newtcolorbox{puentebox}{
  enhanced,breakable,pad at break=2mm,colback=milcGradeSoft,colframe=milcGradeDark,
  borderline west={3pt}{0pt}{milcGradeDark},
  arc=0mm,boxrule=0pt,left=5mm,right=4mm,top=2.5mm,bottom=2.5mm,
  before skip=4pt,after skip=8pt,
  fontupper=\itshape\small\color{milcNegro}\RaggedRight
}
% La flecha va en modo matematico a proposito: Helvetica Neue NO tiene el glifo
% U+2192, asi que \textrightarrow se compilaba a NADA («Missing character: There
% is no → in font Helvetica Neue») y el puente salia sin flecha. En modo
% matematico la flecha viene de la fuente matematica y si se dibuja.
\newcommand{\puente}[1]{%
  \begin{puentebox}%
    \textcolor{milcGradeDark}{$\rightarrow$}\hspace{2mm}#1%
  \end{puentebox}%
}

\newcommand{\linead}[1]{\noindent\color{milcLinea}\rule{#1}{0.5pt}\color{milcNegro}}
\newcommand{\respuesta}[1]{\par\vspace{2mm}\foreach \n in {1,...,#1}{\noindent\color{milcLinea}\rule{\linewidth}{0.5pt}\par\vspace{5mm}}\color{milcNegro}}
\newcommand{\checkbox}{\textcolor{milcGradeDark}{\fbox{\phantom{x}}}\hspace{5pt}}

% ─── Listas aireadas para las guías socioemocionales ────────────────────
% Componer las verificaciones como «\checkbox texto\\» deja el texto que salta
% de línea debajo del cuadrito y apelmaza el bloque. Estos entornos alinean la
% continuación y airean los ítems. Los usa Territorio Interior; cualquier
% familia puede hacerlo.
\newlist{checklista}{itemize}{1}
\setlist[checklista]{
  label=\textcolor{milcGradeDark}{\fbox{\phantom{x}}},
  leftmargin=9mm, labelsep=3.2mm, itemsep=2.4mm,
  topsep=2.5mm, parsep=0pt, partopsep=0pt, before=\RaggedRight
}
\newlist{pasoslista}{enumerate}{1}
\setlist[pasoslista]{
  label=\textcolor{milcGradeDark}{\sffamily\bfseries\arabic*.},
  leftmargin=8mm, labelsep=3mm, itemsep=2.8mm,
  topsep=1mm, parsep=0pt, partopsep=0pt, before=\RaggedRight
}

\newcommand{\phasecard}[4]{
\begin{tcolorbox}[enhanced,colback=#3,colframe=#2,arc=3mm,boxrule=.5pt,height=3.05cm,valign=center,halign=center,left=1.5mm,right=1.5mm,top=2mm,bottom=2mm]
{\sffamily\bfseries\fontsize{22}{24}\selectfont\textcolor{#2}{#1}}\par
{\sffamily\bfseries\small #4}
\end{tcolorbox}
}

% ── Piezas del rediseno 2026 (legibilidad 6.º-11.º) ───────────────────
% Banda de estacion: reemplaza el titlebox macizo. Titulo a la izquierda,
% dato practico (tiempo/modalidad) a la derecha, en una sola linea.
% Nunca huerfana: pide 9 lineas libres (o salta de pagina rellenando con
% \vfill) y prohibe el salto justo despues de la banda. Nueve y no seis:
% la caja partible que sigue necesita ~80pt (titulo adosado, relleno y dos
% lineas) para arrancar en la misma pagina; con menos, tcolorbox la manda
% entera a la siguiente y la banda se queda sola. El penalty 10000 va
% en `after`, ANTES del espacio posterior: un `after skip` (aunque sea 0pt)
% es un \vskip tras la caja, y ese glue es un punto de corte legal que un
% \nopagebreak escrito despues de \end{tcolorbox} ya no cierra. Cada banda
% deja marcador PDF y actualiza la cabecera derecha.
\newcounter{milcestacion}
\newcommand{\milcestacion}[2]{%
\Needspace*{9\baselineskip}%
\stepcounter{milcestacion}%
\pdfbookmark[1]{#2}{milcestacion\themilcestacion}%
\markright{#2}%
\begin{tcolorbox}[enhanced,colback=#1,colframe=#1,arc=2mm,boxrule=0pt,
  left=5mm,right=5mm,top=2.6mm,bottom=2.6mm,before skip=11pt,
  after={\par\nopagebreak\vskip7pt}]
{\sffamily\bfseries\fontsize{15}{18}\selectfont\milctintasobre{#1}#2}
\end{tcolorbox}}

% Tarjeta de paso numerada: sustituye las tablas «Pilar / Aplicacion», que
% eran formato de consulta docente, no de estudiante. El numero va en un
% circulo sobre el borde para que la secuencia se lea de un vistazo.
\newcommand{\milcpaso}[3]{%
\Needspace{4\baselineskip}%
\begin{tcolorbox}[enhanced,colback=white,colframe=milcLinea,arc=1.5mm,boxrule=.9pt,
  left=14mm,right=4mm,top=3mm,bottom=3mm,before skip=4pt,after skip=4pt,
  fontupper=\RaggedRight,
  overlay={\node[anchor=north west,circle,fill=#2,text=white,inner sep=0pt,
    minimum size=7.4mm,font=\sffamily\bfseries\small]
    at ([xshift=3.6mm,yshift=-3.2mm]frame.north west) {#1};}]
#3
\end{tcolorbox}}

% Casilla marcable de tamano util con lapiz (la anterior era \fbox{\phantom{x}},
% de alto variable segun el texto que la seguia).
\newcommand{\milcmarca}{\framebox[3.8mm]{\rule{0pt}{3.4mm}}}

% Fila marcable de lista suelta (checks de la estacion 1).
\newcommand{\milccheckrow}[1]{%
\par\vspace{1.8mm}\noindent
\makebox[6.4mm][l]{\raisebox{-.55mm}{\textcolor{milcNegro}{\milcmarca}}}%
\parbox[t]{\dimexpr\linewidth-6.4mm\relax}{\RaggedRight #1}\par}

% Celda marcable dentro de la rubrica (tabla de autoevaluacion). Misma
% sangria francesa, pero pensada para vivir dentro de una columna X.
\newcommand{\milcopcion}[1]{%
\makebox[5.2mm][l]{\raisebox{-.55mm}{\textcolor{milcNegro}{\milcmarca}}}%
\parbox[t]{\dimexpr\linewidth-5.2mm\relax}{\RaggedRight #1}}

% ── Recursos visuales de la guía ──────────────────────────────────────
% Declarados en el bloque `recursos:` del YAML y emitidos por el builder.
% \guiaFigura   → imagen rasterizada o PDF (foto, carta celeste, montaje).
% \guiaDiagrama → fuente TikZ/circuitikz (.tex) incrustada como vector:
%                 es la vía para circuitos y esquemáticos editables.
% [#1] = texto alternativo (alt) · #2 = ruta relativa al .tex · #3 = pie
% El alt llega al PDF etiquetado (/Alt) y lo lee el lector de pantalla.
\newcommand{\guiaFigura}[3][]{%
\begin{center}
\includegraphics[width=0.88\linewidth,height=0.38\textheight,keepaspectratio,alt={#1}]{#2}\\[1.5mm]
{\footnotesize\itshape\textcolor{milcNegro!75}{#3}}
\end{center}
}

\newcommand{\guiaDiagrama}[2]{%
\begin{center}
\input{#1}\\[1.5mm]
{\footnotesize\itshape\textcolor{milcNegro!75}{#2}}
\end{center}
}

\begin{document}
\thispagestyle{empty}

% ── Portada ───────────────────────────────────────────────────────────
% Antes: un rectangulo de color a pagina completa. Se veia bien en pantalla,
% pero en la fotocopiadora del colegio se comia el toner y salia gris sucio.
% Ahora la banda ocupa el tercio superior y el resto va en blanco.
\begin{tikzpicture}[remember picture,overlay]
  \path[fill=milcGradePrimary]
    (current page.north west) rectangle ([yshift=-10.6cm]current page.north east);

  \node[anchor=north west,text=milcGradeText] at ([xshift=2.0cm,yshift=-1.55cm]current page.north west)
    {\sffamily\bfseries\fontsize{12}{14}\selectfont GRADO};
  \node[anchor=north west,text=milcGradeText,opacity=.85] at ([xshift=2.0cm,yshift=-2.35cm]current page.north west)
    {\sffamily\bfseries\fontsize{8.5}{10}\selectfont SERIE GUÍAS · TECNOLOGÍA E INFORMÁTICA};
  \node[anchor=north east,text=milcGradeText,opacity=.85,align=right] at ([xshift=-2.0cm,yshift=-1.55cm]current page.north east)
    {\sffamily\bfseries\fontsize{9}{12}\selectfont I.E. Sor María Juliana\\Cartago, Valle del Cauca};

  \node[anchor=west,text=milcGradeText] (gradonum) at ([xshift=1.85cm,yshift=-7.1cm]current page.north west)
    {\sffamily\bfseries\fontsize{112}{112}\selectfont 9};
  % El circulito de grado (°) solo aplica a las guias de grado («11°»); en
  % Bebras y Semillero el numero grande es un momento/modulo, no un grado.
  % Se posiciona relativo al nodo `gradonum`, no en coordenadas absolutas.
  \draw[milcGradeText,line width=1.6pt]
    ([xshift=2.0mm,yshift=-4.5mm]gradonum.north east) circle (.15cm);

  \node[anchor=west,text=milcGradeText,text width=11.4cm,align=left] at ([xshift=1.5cm,yshift=0.5cm]gradonum.east)
    {
     {\sffamily\bfseries\fontsize{9}{11}\selectfont\MakeUppercase{Período 3 · Datos --- del registro al insight}}\\[3mm]
     {\sffamily\bfseries\fontsize{21}{25}\selectfont\setlength{\baselineskip}{27pt}Gráficos ---\\[4pt]la imagen\\[4pt]de los datos}\\[3.5mm]
     {\sffamily\bfseries\fontsize{15}{18}\selectfont Guía 26-9-TIC}
    };
\end{tikzpicture}

\vspace*{9.4cm}
\vfill

{\sffamily\bfseries\fontsize{8}{10}\selectfont\textcolor{milcGradeDark}{LO QUE TE LLEVAS HOY}}

\begin{tcolorbox}[enhanced,colback=white,colframe=milcNegro,arc=2mm,boxrule=1.6pt,
  left=5mm,right=5mm,top=4mm,bottom=4mm,before skip=3pt,after skip=12pt,
  fontupper=\RaggedRight\fontsize{11}{15.5}\selectfont]
3 gráficos (barra, línea, circular) generados a partir de la tabla del estudiante, cada uno acompañado de la justificación de por qué ese tipo de gráfico era el adecuado para esa pregunta
\end{tcolorbox}

\begin{tcolorbox}[enhanced,colback=milcGris,colframe=milcGris,arc=2mm,boxrule=0pt,
  left=5mm,right=5mm,top=3.5mm,bottom=3.5mm,after skip=12pt]
{\sffamily\bfseries\fontsize{8}{10}\selectfont\textcolor{milcNegro!55}{ESTA GUÍA ES DE}}\\[3mm]
\begin{tabularx}{\linewidth}{X p{3.2cm} p{3.2cm}}
\linead{\linewidth} & \linead{3.0cm} & \linead{3.0cm}\\[-1mm]
{\fontsize{8.5}{10}\selectfont\textcolor{milcNegro!55}{Nombre completo}} &
{\fontsize{8.5}{10}\selectfont\textcolor{milcNegro!55}{Grupo}} &
{\fontsize{8.5}{10}\selectfont\textcolor{milcNegro!55}{Fecha}}\\
\end{tabularx}
\end{tcolorbox}

\vfill

\noindent\textcolor{milcLinea}{\rule{\linewidth}{1pt}}\\[2mm]
{\sffamily\bfseries\fontsize{9.5}{12}\selectfont Autor: Dr. Álvaro Cárdenas Orozco}\\[1mm]
{\sffamily\fontsize{8.5}{11}\selectfont\textcolor{milcNegro!55}{Metodología MILC · apertura ancestral · pensamiento computacional · triángulo Dussel–estoicismo–Floridi}}

\newpage

\section*{Apertura: Gráficos --- la imagen de los datos}

\begin{softbox}{milcGradeDark}{milcGradeSoft}{Saber ancestral}
Antes del gráfico de barras existieron los \textbf{petroglifos Calima} del norte del Valle del Cauca (siglos II a XII) que mostraban caminos, ríos y zonas de cosecha en piedra; los \textbf{dibujos en arena} de los pescadores afro-pacíficos del Chocó y Buenaventura para explicar corrientes y zonas de pesca; los \textbf{esquemas de la abuela campesina de Cartago} dibujados en el suelo con un palo para mostrar el patrón del año productivo; los \textbf{tejidos Wayuu} cuyos rombos y triángulos cuentan historias de tribu; la \textbf{balsa de oro Muisca} del altiplano cundiboyacense que es ella misma un mapa simbólico de jerarquía y ritual. Visualizar datos es un gesto humano antiguo. El software no inventó la gráfica: cambió el medio y la velocidad, pero la pregunta sigue siendo la misma --- \emph{¿cómo le muestro a otro lo que yo ya vi en los números?}
\end{softbox}

\begin{softbox}{milcTurquesa}{milcGris}{Saber contemporáneo}
Un \textbf{gráfico} convierte una tabla en una imagen para que el ojo lea de un vistazo lo que la tabla solo dice fila por fila. Hay 3 tipos básicos que resuelven el 80~\% de los casos: (1) \texttt{Barras} para \emph{comparar} categorías; (2) \texttt{Líneas} para mostrar \emph{evolución} en el tiempo; (3) \texttt{Circular} (\emph{pie}) para mostrar \emph{participación} sobre un total. Elegir mal el tipo de gráfico no es un detalle estético: hace que la pregunta se responda mal o que el lector entienda otra cosa.
\end{softbox}

\begin{softbox}{milcMostaza}{milcGris}{Pregunta-puente}
\textit{¿Cómo hacía el pescador afro-pacífico para enseñar la corriente con un solo dibujo en la arena? ¿Y qué pierde un analista novato cuando hace un gráfico bonito sin haber escrito antes la pregunta que el gráfico debe responder?}
\end{softbox}

\begin{softbox}{milcVerde}{milcGris}{Saber y saber hacer}
Al terminar podrás: (1) \textbf{identificar} los 3 tipos básicos de gráfico (barra, línea, circular) y la pregunta que responde cada uno; (2) \textbf{analizar} cuándo un gráfico engaña al ojo y cuándo lo orienta bien; (3) \textbf{aplicar} la construcción de los 3 gráficos sobre tu propia tabla; (4) \textbf{crear} 3 gráficos con justificación escrita de por qué el tipo elegido era el adecuado.
\end{softbox}



\small
\begin{tabularx}{\textwidth}{>{\bfseries}p{3.0cm}Y}
\rowcolor{milcGradePrimary!12}Autor & Dr. Álvaro Cárdenas Orozco\\
DBA & Tratamiento y visualización honesta de datos para la toma de decisiones (MEN, T\&I)\\
Referentes & Ética del dato · Visualización honesta · Pensamiento computacional\\
Producto final & 3 gráficos (barra, línea, circular) generados a partir de la tabla del estudiante, cada uno acompañado de la justificación de por qué ese tipo de gráfico era el adecuado para esa pregunta\\
\end{tabularx}
\normalsize

\puente{Antes de empezar, mira el viaje completo. La sesión 5 te enseñó a agrupar la tabla en una tabla dinámica; hoy aprendes a \emph{ponerle imagen} a esa agrupación. Lo mismo que viste en filas y columnas, ahora lo vas a ver en altura de barra, pendiente de línea y porción de torta. Es el salto del \emph{número} a la \emph{imagen}.}

\milcestacion{milcGradeDark}{Ruta MILC: así vamos a aprender}
\begin{tabularx}{\textwidth}{>{\centering\arraybackslash}X>{\centering\arraybackslash}X>{\centering\arraybackslash}X>{\centering\arraybackslash}X}
\phasecard{1}{milcVerde}{milcGris}{Escucha\\[-1mm]\small Observo, escucho y pregunto.} &
\phasecard{2}{milcTurquesa}{milcGris}{Entender\\[-1mm]\small Sistematizo y ordeno.} &
\phasecard{3}{milcMagenta}{milcGris}{Hacer\\[-1mm]\small Construyo y aplico.} &
\phasecard{4}{milcVino}{milcGris}{Cierre\\[-1mm]\small Evalúo y reflexiono.}
\end{tabularx}


\puente{Antes de teorizar, vas a mirar 3 gráficos (uno bueno y dos engañosos) y a decir qué entendiste de cada uno. Después vas a ver la tabla detrás de ellos y a comparar. Ese contraste te muestra que un gráfico puede iluminar o engañar.}

\milcestacion{milcVerde}{Estación 1 de 3 · Escucha}

\begin{softbox}{milcVerde}{milcGris}{Escena para escuchar}
\textbf{Actividad 1 · IDENTIFICA --- ¿Qué entendiste de este gráfico?} (15 min · individual). El docente muestra 3 gráficos: uno bien hecho (barras con eje completo), uno con eje Y cortado que exagera diferencias, y uno circular con 12 porciones donde no se entiende nada. Sin ver la tabla detrás, responde para cada gráfico: \emph{¿qué pregunta responde?}, \emph{¿qué conclusión sacarías?}, \emph{¿qué te parece sospechoso o confuso?}. Después se muestran las tablas y comparas tu lectura con los datos reales.
\end{softbox}

{\sffamily\bfseries\fontsize{8}{10}\selectfont\textcolor{milcNegro!55}{MARCA CUANDO LO HAYAS HECHO}}
\milccheckrow{Anoté qué pregunta parecía responder cada uno de los 3 gráficos.}
\milccheckrow{Anoté qué conclusión sacaría de cada uno antes de ver la tabla.}
\milccheckrow{Identifiqué cuál gráfico era confiable y cuáles me engañaron, y por qué.}

\begin{infoband}{milcVerde}


\textbf{Lo que va al cuaderno (Actividad 1 · IDENTIFICA).} Título: «Actividad 1 --- ¿Qué entendiste de este gráfico?». \textbf{Formato:} tabla de 3 filas (un gráfico por fila) y 3 columnas (Pregunta que parecía responder~|~Conclusión a primera vista~|~Qué cambió cuando vi la tabla). \textbf{Sabes que terminaste cuando} reconoces que un mismo dato puede contar dos historias distintas según el tipo de gráfico.
\end{infoband}

\puente{Ya viste el contraste. Ahora vas a entender la \textbf{anatomía} de un gráfico bien hecho: qué pregunta responde, qué tipo elige, qué títulos lleva, qué ejes muestra, qué leyenda usa, qué cosas \textbf{no} debe hacer (como cortar el eje Y o usar circular con 15 categorías).}

\milcestacion{milcTurquesa}{Estación 2 de 3 · Entender}

\begin{softbox}{milcTurquesa}{milcGris}{Lo que vas a entender}
Los 3 tipos básicos de gráfico responden preguntas distintas. Vamos a descomponer cada uno con los pilares del pensamiento computacional para que sepas exactamente cuándo elegir cuál.
\end{softbox}

\milcpaso{1}{milcTurquesa}{Los 3 tipos básicos se descomponen así: (1) \textbf{Barras} responde \emph{¿cuál categoría es mayor o menor?} --- eje X = categorías, eje Y = valor; ideal cuando las categorías no son tiempo. (2) \textbf{Líneas} responde \emph{¿cómo cambia en el tiempo?} --- eje X = fecha o periodo ordenado, eje Y = valor; ideal para evolución. (3) \textbf{Circular} responde \emph{¿qué porción del total es cada parte?} --- una sola categoría con pocas porciones (idealmente 3-5, máximo 7); ideal para participación.}
\milcpaso{2}{milcTurquesa}{Todo gráfico sigue el patrón \textbf{``pregunta $\to$ tipo $\to$ ejes $\to$ títulos''}: primero defines la pregunta, después eliges el tipo, después configuras los ejes (con cero en el eje Y para barras y líneas, salvo justificación), y por último escribes un título que diga qué muestra el gráfico, no qué tipo es. \emph{``Gasto mensual en transporte (enero-mayo)''} es título; \emph{``Gráfico de barras 1''} no lo es.}
\milcpaso{3}{milcTurquesa}{Lo esencial cabe en una regla: \emph{el tipo de gráfico debe coincidir con la pregunta}. Comparar 5 cursos: barras. Ver evolución de notas durante el periodo: líneas. Mostrar la participación de 4 áreas en el total de horas: circular. Si forzás un circular con 15 categorías, el lector no entiende nada; si forzás una línea para comparar 4 productos sin tiempo, conectas puntos que no tienen relación temporal.}
\milcpaso{4}{milcTurquesa}{Algoritmo: (1) escribe la pregunta; (2) decide si comparas (barras), si miras evolución (líneas) o si muestras participación (circular); (3) selecciona la tabla o la tabla dinámica; (4) Insertar $\to$ Gráfico recomendado; (5) ajusta título, ejes y leyenda; (6) verifica que el eje Y arranque en cero salvo motivo declarado; (7) escribe la justificación.}

\begin{drawbox}{milcTurquesa}{Anatomía de un gráfico bien hecho}
\textbf{Pregunta:} qué responde el gráfico en 1 frase. \\[1mm] \textbf{Tipo:} barra / línea / circular. \\[1mm] \textbf{Ejes:} eje X y eje Y nombrados con unidad. \\[1mm] \textbf{Eje Y:} arranca en cero (salvo justificación visible). \\[1mm] \textbf{Título:} dice qué muestra, no qué tipo es. \\[1mm] \textbf{Leyenda:} corta y legible; solo si hay más de 1 serie. \\[1mm] \textbf{Justificación:} 1 frase con por qué este tipo y no otro.
\end{drawbox}

\begin{softbox}{milcMostaza}{milcGris}{Errores comunes y su corrección}
(1) Cortar el eje Y para exagerar diferencias: distorsiona la lectura. (2) Usar circular con más de 7 porciones: ya nadie distingue las porciones. (3) Conectar con líneas datos que no son temporales: sugiere evolución inexistente. (4) Olvidar título y unidades en los ejes: el lector no sabe qué está mirando. (5) Hacer un gráfico bonito sin pregunta escrita: decoración, no comunicación.
\end{softbox}

\begin{infoband}{milcTurquesa}


\textbf{Lo que va al cuaderno (Actividad 2 · EVALÚA).} Título: «Actividad 2 --- Cuál tipo de gráfico engaña en cada caso». \textbf{Formato:} 3 fichas, una por tipo (barra, línea, circular), cada una con: pregunta que responde + cuándo elegirlo + cuándo NO elegirlo + 1 error frecuente. \textbf{Veredicto (3 casos):} para estos 3 datos, decide qué gráfico usar y cuál NO usar (con justificación de 1 razón cada uno): (1) ventas mensuales del año en una tienda; (2) porcentaje de tipos de sangre en Colombia; (3) altura promedio por edad de un estudiante de 6° a 11°. \textbf{Sabes que terminaste cuando} puedes elegir el tipo correcto sin consultar y \emph{explicar por qué los demás engañarían}.
\end{infoband}

\puente{Tienes el método. Toca aplicarlo: 3 gráficos (barra, línea, circular) sobre tu propia tabla o sobre la tabla dinámica de Sesión 5, cada uno con justificación escrita de por qué elegiste ese tipo y no otro.}

\milcestacion{milcMagenta}{Estación 3 de 3 · Hacer}

\begin{softbox}{milcMagenta}{milcGris}{Cómo vas a construir tu producto}
\textbf{Actividad 3 · CREA --- 3 gráficos con justificación} (30 min · individual). Hora de crear. Sobre tu tabla de Sesión 2 o la tabla dinámica de Sesión 5, vas a generar 3 gráficos --- uno de cada tipo --- cada uno con su pregunta clara y su justificación escrita de por qué elegiste ese tipo y no otro.
\end{softbox}

\milcpaso{1}{milcMagenta}{Tu trabajo se descompone en 3 piezas, una por gráfico: (1) un gráfico de barras que compara categorías; (2) un gráfico de líneas que muestra evolución temporal; (3) un gráfico circular que muestra participación sobre un total. Variar los 3 tipos te obliga a entender por qué cada uno existe.}
\milcpaso{2}{milcMagenta}{Sigue el patrón \textbf{``pregunta antes que gráfico''}: para cada uno, escribe primero la pregunta en lenguaje natural. \emph{``¿Cuál es el área con más horas asignadas?''} --- barras. \emph{``¿Cómo evolucionó mi promedio durante el periodo?''} --- líneas. \emph{``¿Qué porcentaje del gasto fue en transporte?''} --- circular. Sin pregunta, no hay gráfico.}
\milcpaso{3}{milcMagenta}{Cada justificación expresa una sola idea, no tres. \emph{``Elegí barras porque comparo categorías que no son tiempo''} ya es suficiente. Si justificás con párrafo largo, recortá. La claridad técnica vence al adorno argumental.}
\milcpaso{4}{milcMagenta}{Algoritmo: (1) escribe pregunta; (2) elige tipo; (3) selecciona tabla; (4) inserta gráfico; (5) configura título y ejes; (6) revisa que el eje Y arranque en cero; (7) toma captura; (8) escribe justificación de 1-2 frases. Repite 3 veces.}

\begin{drawbox}{milcMagenta}{Checklist antes de cerrar los 3 gráficos}
\checkbox 3 preguntas escritas, una por gráfico, antes de insertar. \\[1mm] \checkbox 3 gráficos creados, uno de cada tipo (barra, línea, circular). \\[1mm] \checkbox Todos con título descriptivo y ejes con unidad. \\[1mm] \checkbox Eje Y arranca en cero en barras y líneas (o justificación declarada). \\[1mm] \checkbox 3 justificaciones escritas que digan por qué ese tipo y no otro.
\end{drawbox}

\begin{softbox}{milcMagenta}{milcGris}{Plantilla de trabajo}
\textbf{Gráfico 1 (barras).} \textit{Pregunta + categorías comparadas + justificación: ``Elegí barras porque...''.} \\[1mm] \textbf{Gráfico 2 (líneas).} \textit{Pregunta + periodo evaluado + justificación: ``Elegí líneas porque hay evolución temporal y...''.} \\[1mm] \textbf{Gráfico 3 (circular).} \textit{Pregunta + total considerado + justificación: ``Elegí circular porque muestro participación sobre un total y son pocas porciones...''.}
\end{softbox}

\begin{infoband}{milcMagenta}


\textbf{Lo que va al cuaderno (Actividad 3 · CREA).} Título: «Actividad 3 --- Mis 3 gráficos con justificación». \textbf{Formato:} 3 fichas con pregunta + tipo + captura pegada o descrita + justificación de 1-2 frases. \textbf{Extensión:} 1 página de cuaderno. \textbf{Sabes que terminaste cuando} puedes defender frente a un compañero por qué cada gráfico es del tipo correcto y no otro.
\end{infoband}

\puente{Lo que sigue resume \emph{qué} entregas hoy: 3 gráficos + 3 justificaciones. El criterio principal: que la justificación nombre \textbf{la pregunta} que responde el gráfico y \textbf{por qué otro tipo de gráfico habría sido peor}.}

\milcestacion{milcMostaza}{Producto de la sesión}
\textbf{3 gráficos (barra, línea, circular) con justificación escrita}.
\begin{softbox}{milcMostaza}{milcGris}{Criterios de calidad}
(1) 3 preguntas escritas antes de crear los gráficos. (2) 3 gráficos, uno de cada tipo, sobre la tabla propia o la tabla dinámica. (3) Títulos descriptivos y ejes con unidad en cada gráfico. (4) Eje Y arranca en cero (o justificación declarada). (5) 3 justificaciones que nombren por qué ese tipo y no otro.
\end{softbox}

\puente{Antes de cerrar, mira los gráficos desde las cinco dimensiones humanas. Graficar bien es respeto al lector; graficar mal o con trampas es ruido visual o manipulación.}

\milcestacion{milcVino}{Evaluación liberadora · cinco dimensiones}

\begin{softbox}{milcVino}{milcGris}{Sentido de la evaluación}
Evaluar no es solo revisar una nota. Es reconocer qué aprendí, qué cambió en mí, qué emociones manejé, cómo me relaciono con la ciudad y la comunidad, y qué le dejaré a quien viene detrás.
\end{softbox}

\textbf{Mi reflexión liberadora en cinco dimensiones.} Completa cada frase iniciada:

\small
\begin{tabularx}{\textwidth}{>{\bfseries\RaggedRight\arraybackslash}p{3.4cm}Y>{\RaggedRight\arraybackslash}p{4.6cm}}
\rowcolor{milcVino}\color{white}Dimensión & \color{white}\bfseries Mi respuesta (frame starter) & \color{white}\bfseries Algo que descubrí\\
1. Personal & Lo que más cambió en mí fue \linead{6.5cm} & \linead{4.2cm}\\
2. Emocional & La emoción más fuerte fue \linead{2.5cm} y la regulé \linead{3cm} & \linead{4.2cm}\\
3. Ciudadana & En democracia esto importa porque \linead{6cm} & \linead{4.2cm}\\
4. Local (Cartago) & En mi barrio sería útil para \linead{6.5cm} & \linead{4.2cm}\\
5. Intergeneracional & A mi abuela/abuelo le diría \linead{6.5cm} & \linead{4.2cm}\\
\end{tabularx}
\normalsize

\begin{infoband}{milcVino}
\textbf{Para la próxima sustentación.} Vuelve a esta tabla en seis meses. Verás que las respuestas cambian — no porque lo aprendido cambie, sino porque tú cambias. La evaluación liberadora es una conversación contigo a través del tiempo.
\end{infoband}


\puente{Y para terminar, tres voces sobre la imagen de los datos. Dussel sobre los gráficos que invisibilizan, Séneca sobre la sencillez como virtud visual, Floridi sobre la imagen como infraestructura de la opinión pública contemporánea.}

\milcestacion{milcGradeDark}{Triángulo de pensamiento · cierre filosófico}

\begin{softbox}{milcVino}{milcGris}{Dussel · lente del nosotros}
\textit{Toda imagen muestra y oculta al mismo tiempo: la pregunta es qué oculta.}

{\footnotesize\itshape\color{milcNegro!70}--- formulación de esta guía, al modo de Enrique Dussel. No es una cita textual.}

Cuando haces un gráfico circular con \emph{``otros''} agrupando 8 categorías pequeñas, esas 8 categorías dejan de existir para el lector. Cuando un noticiero usa un eje Y cortado para mostrar la \emph{``caída brutal del empleo''}, la imagen exagera lo que la tabla no diría. Graficar es decidir qué entra al ojo y qué queda fuera. La phronesis del analista exige declarar qué simplificación hizo el gráfico y qué se perdió en el camino.

\textbf{En mis 3 gráficos, ¿qué categorías cayeron en \emph{``otros''}? ¿Qué decisión cambiaría si esas categorías fueran visibles una por una?}
\end{softbox}
\linead{\linewidth}\\[1mm]
\linead{\linewidth}

\begin{softbox}{milcMostaza}{milcGris}{Estoicismo · Séneca · cuidado interior}
\textit{Lo que no necesita adorno es lo que dice la verdad.}

{\footnotesize\itshape\color{milcNegro!70}--- formulación de esta guía, al modo de Séneca. No es una cita textual.}

Es tentador llenar el gráfico de colores, sombras, efectos 3D y leyendas largas. La sabiduría estoica recomienda lo contrario: \emph{cuanta menos decoración, más datos se ven}. Un gráfico con título claro, ejes con unidad y una sola serie comunica mejor que un \emph{infográfico} cargado de iconos. Quien grafica sin adornos respeta el tiempo del lector.

\textbf{¿Cuánto adorno habrá en mis gráficos que no ayuda al lector a responder la pregunta? ¿Qué pasaría si lo quito?}
\end{softbox}
\linead{\linewidth}\\[1mm]
\linead{\linewidth}

\begin{softbox}{milcTurquesa}{milcGris}{Floridi · lente de la infoesfera}
\textit{Los gráficos son la infraestructura visual con que las sociedades contemporáneas se forman opinión.}

{\footnotesize\itshape\color{milcNegro!70}--- formulación de esta guía, al modo de Luciano Floridi. No es una cita textual.}

Casi toda noticia económica, política o de salud llega al ciudadano a través de un gráfico. Quien no sabe leerlos críticamente recibe la conclusión del editor, no la información de la tabla. Dominar gráficos en tu propia tabla es el primer paso para criticar (o defender con argumento) los gráficos que te muestran en redes, prensa y publicidad.

\textbf{¿Cuántos gráficos vi esta semana en redes o noticias? ¿Cuántos eran del tipo adecuado para la pregunta que decían responder?}
\end{softbox}
\linead{\linewidth}\\[1mm]
\linead{\linewidth}

\begin{infoband}{milcNegro}
\textbf{Nota para el docente.} Las tres formulaciones son del autor de la guía, no citas textuales de los pensadores: recogen su lente para pensar el tema, no sus palabras. Si vas a citarlos en clase, remite a la obra original.
\end{infoband}

\milcestacion{milcVino}{Mi autoevaluación · marca una casilla por fila}

{\small
\setlength{\extrarowheight}{2.2pt}
\begin{tabularx}{\textwidth}{>{\RaggedRight\arraybackslash\bfseries}p{3.0cm}YYY}
\rowcolor{milcVino}\color{white}\bfseries Criterio & \color{white}\bfseries Lo logré & \color{white}\bfseries En proceso & \color{white}\bfseries Necesito apoyo\\[.6mm]
Comprensión del tema & \milcopcion{Explico las ideas con mis palabras.} & \milcopcion{Reconozco las ideas, falta precisión.} & \milcopcion{Necesito repasar lo básico.}\\[.8mm]
\rowcolor{milcGris}Pensamiento computacional & \milcopcion{Usé los pilares al armar mi producto.} & \milcopcion{Usé algunos pilares.} & \milcopcion{Necesito releer la estación 2.}\\[.8mm]
Aplicación y producto & \milcopcion{Mi producto cumple los criterios.} & \milcopcion{Está armado, con ajustes pendientes.} & \milcopcion{Aún está en construcción.}\\[.8mm]
\rowcolor{milcGris}Reflexión 5 dimensiones & \milcopcion{Respondí cada una con honestidad.} & \milcopcion{Respondí algunas con detalle.} & \milcopcion{Mis respuestas son muy generales.}\\[.8mm]
Triángulo de pensamiento & \milcopcion{Conecté las 3 voces con mi trabajo.} & \milcopcion{Conecté al menos una.} & \milcopcion{No las relaciono con mi proceso.}\\[.8mm]
\end{tabularx}}

\vspace{3mm}
\textbf{Hoy entendí que:}\\[1mm]
\linead{\linewidth}\\[1mm]
\linead{\linewidth}

\textbf{Mi compromiso para la próxima sesión es:}\\[1mm]
\linead{\linewidth}\\[1mm]
\linead{\linewidth}

\begin{softbox}{milcMostaza}{milcGris}{Cita de despedida · Marco Aurelio}
\textit{``El obstáculo en el camino se vuelve el camino. Lo que estorba al camino se vuelve camino.''}\\[1mm]
Cada error de hoy, bien observado, se convierte en pieza del camino que pisarás mañana.
\end{softbox}

\small
\begin{tabularx}{\textwidth}{>{\bfseries}p{3.6cm}Y}
\rowcolor{milcGradeDark!12}Nombre completo: & \linead{10cm}\\
Fecha: & \linead{4cm} \quad Curso/grupo: \linead{4cm}\\
Mi firma: & \linead{9cm}\\
\end{tabularx}
\normalsize

\begin{infoband}{milcGradeDark}
\textbf{Para la próxima semana.} Toma un gráfico real que veas en prensa, redes o un cuaderno escolar. Identifica: \emph{¿qué pregunta dice responder?}, \emph{¿es el tipo adecuado?}, \emph{¿el eje Y arranca en cero?}, \emph{¿qué categoría queda escondida en ``otros''?}. La phronesis no se queda graficando: se entrena leyendo los gráficos que rigen las decisiones del día.
\end{infoband}

\end{document}
